% ===== sec/06_discussion.tex =====
\section{Discussion and Limitations}
\label{sec:discussion}

\subsection{Limitations}

Three limitations qualify the results.  First, the Type-0 abstraction
for $\mathcal{G}$ is an idealisation: the concrete laptop execution
is bounded and hence formally an LBA.  Lemma~\ref{lem:bounded-soundness}
supports the one-way comparison but does not prove strict Type-0
membership.  Second, the architecture has been instantiated on one
task (\textsc{Spaceship Titanic}); generalisation to other domains
(images, text, open-ended code) may require an expanded start
language and keyword set.  Third, the controlled ablation uses
synthetic tasks and simulated judges rather than the full
\textsc{MLE-bench} harness with live API judges.

\subsection{Threats to validity}

\paragraph{Construct validity.}
The trace-language identity criterion (Definition~\ref{def:traceeq})
discards information about chain-of-thought and intermediate JSON.
We accept the discard because only the trace is verifiable.

\paragraph{Internal validity.}
The LLM core is non-deterministic; leaderboard placement varies by
$\pm$10--20 ranks across re-runs.  We report a single placement
($710/2330$) and acknowledge the variance.  The reproducibility kit
allows external re-runs at fixed random seeds.

\paragraph{External validity.}
We expect the trace-language account to generalise to other
Kaggle-style tabular tasks because the start language and verifier
are task-agnostic.  Generalisation outside the tabular regime is not
claimed.

\subsection{Threats to theoretical validity}

\paragraph{Single trace problem.}
An empirical trace language is always a finite set; the formal
language it belongs to is infinite.  Observed traces do not
determine the Chomsky class of the generating language: a finite
sample is consistent with every level of the hierarchy.
Consequently, the class assignments in
Table~\ref{tab:class-signature} are modelling hypotheses, not
membership proofs.  We make this explicit by appending ``?'' to
every non-Type-3 class label.  The verifier is the sole exception
because its recogniser (a concrete DFA) is explicit.

\paragraph{Abstraction gap.}
The Type-0 abstraction for $\mathcal{G}$ is an analytical device:
it assumes unbounded tape and time, whereas the concrete system
runs on a finite laptop with a fixed timeout.  Lemma~\ref{lem:bounded-soundness}
establishes one-way soundness (bounded $\subseteq$ idealised) but
not the converse.  A trace that the idealised machine could produce
after $10^{10^{10}}$ steps may be practically unreachable.
The abstraction is useful precisely because verification claims
discharged against the idealised machine apply a fortiori to the
bounded one; but the abstraction cannot be used to argue that
$\mathcal{G}$ is Turing-complete in the concrete sense.

\paragraph{Projection ambiguity.}
Different generators may induce identical projected traces on a
given sub-alphabet $\Sigma'$ while differing arbitrarily on the
complement.  The projection operation
(Definition~\ref{def:projection}) therefore loses information:
two agents that produce the same planner trace may have
qualitatively different researcher traces.  The framework does not
resolve this ambiguity; it records it as a modelling uncertainty
that can be narrowed only by expanding the keyword set of the
verifier or the start alphabet.

\subsection{Future work}

One observed repair loop in \texttt{sop\_orchestrator} is suggestive
of indexed grammars~\cite{aho1968indexed}; we leave strict indexed
membership for future work.  Additional directions include repeating
the ablation on live \textsc{MLE-bench} tasks, studying which start
languages produce stable bindings, and characterising the class of
regular keyword-update rules that preserve decidability under a
Type-3 producer.

\subsection{Closing remark}

Agent behaviour can be represented as a trace language, and a
regular verifier can constrain a more expressive generator while
remaining computationally tractable.  This viewpoint gives useful
producer--consumer tools for language-constrained trace validation.
